\chapter{Décimos y centésimos}\label{ch:g4:decimals}

Entre $3$ y $4$ los números enteros están en silencio --- pero un corredor puede
ser $3$ metros y un poco más adelante. Cortando la unidad en diez, luego una
cien, da el \emph{decimal fracciones}, y una maravillosa
abreviatura para ellos: el \omterm{def:g4:decimals:point}{coma decimal}. (Los números decimales obtienen su
capítulo completo en \cref{ch:g5:decimals}.)

\section{fracciones decimales}

\begin{definition}[Décimas y centésimas]\label{def:g4:decimals:def}
Cortar una unidad en $10$ partes iguales dan \emph{décimas},
$\frac{1}{10}$; cortando cada décimo en $10$ otra vez da
\emph{centésimas}, $\frac{1}{100}$. Entonces
\[
1 = \frac{10}{10} = \frac{100}{100},
\qquad
\frac{1}{10} = \frac{10}{100}.
\]
\end{definition}

\begin{omfigure}
\begin{tikzpicture}[scale=0.42]
  \draw[gray!50, thin] (0,0) grid (10,10);
  \foreach \i in {0,1,2,3,4,5,6,7,8,9} \fill[omDef!35] (3,\i+0.06)
    rectangle (3.94,\i+0.94);
  \fill[omThm!70] (6.05,4.05) rectangle (6.95,4.95);
  \draw[thick] (0,0) rectangle (10,10);
  \node[below, font=\small] at (5,-0.4)
    {the whole square $= 1$;\ \ a column (blue) $= \frac{1}{10}$;\ \
    a small cell (red) $= \frac{1}{100}$};
\end{tikzpicture}

{\small La centena cuadrada: diez columnas de diez celdas. Una columna es una
décimo, una celda por centésima.}
\end{omfigure}

\begin{example}\label{ex:g4:decimals:count}
Sombreado $3$ columnas y $7$ celdas adicionales de las sombras de cien \omterm{def:g3:shapes:shapes}{cuadrados}
\[
\frac{3}{10} + \frac{7}{100} = \frac{30}{100} + \frac{7}{100}
= \frac{37}{100}
\]
del cuadrado --- treinta y siete centésimas.
\end{example}

\section{El punto decimal}

\begin{definition}[escritura decimal]\label{def:g4:decimals:point}
el número $2 + \frac{3}{10} + \frac{7}{100}$ esta escrito
\[
2.37
\]
--- el \emph{coma decimal}\index{coma decimal} separa toda la parte ($2$) de la
parte decimal; el primer dígito después del punto cuenta las décimas, el
segundo las centésimas. Lectura: ``dos coma tres siete'', o ``dos
y treinta y siete centésimas''.
\end{definition}

\begin{example}\label{ex:g4:decimals:convert}
\[
\frac{7}{10} = 0.7, \qquad
\frac{43}{100} = 0.43, \qquad
\frac{5}{100} = 0.05, \qquad
3 + \frac{4}{10} = 3.4 .
\]
Cuidado con $\frac{5}{100}$: cinco \emph{centésimas} necesito un $0$ en
el lugar de las décimas --- $0.05$, no $0.5$.
\end{example}

\begin{omfigure}
\begin{tikzpicture}[scale=1.35]
  \draw[-Stealth] (-0.2,0) -- (8.4,0);
  \foreach \i in {0,10,20,30,40,50,60,70,80}
    \draw ({\i*0.1},-0.09) -- ({\i*0.1},0.09);
  \foreach \i in {0,...,80} \draw ({\i*0.1},-0.045) -- ({\i*0.1},0.045);
  \foreach \x/\l in {0/2, 1/{2.1}, 2/{2.2}, 3/{2.3}, 4/{2.4},
                     5/{2.5}, 6/{2.6}, 7/{2.7}, 8/{2.8}}
    \node[below, font=\footnotesize] at (\x,-0.1) {$\l $};
  \fill[omThm] (3.7,0) circle (1.7pt)
    node[above, font=\small] {$2.37$};
\end{tikzpicture}

{\small Zoom entre $2$ y $2.8$: marcas grandes cada décimo, pequeñas
marca cada centésima. el número $2.37$ se sienta entre $2.3$ y
$2.4$, en el $7$ª marca pequeña.}
\end{omfigure}

\begin{example}[El dinero son centésimas.]\label{ex:g4:decimals:money}
Un céntimo es una centésima de euro (o dólar, o libra): un precio
de $4.85$ \omterm{def:g3:division:half}{medio} $4$ monedas enteras y $85$ centésimas. El dinero es el
\omterm{def:g4:decimals:point}{coma decimal} ya usas todos los días: $10$ centavos $= 0.10$; dos
cincuenta euros $= 2.50$; y $0.05$ ¡Es una moneda de cinco céntimos, no de cincuenta!
\end{example}

\section{Comparando con décimas}

\begin{method}[Comparando decimales (primer contacto)]\label{met:g4:decimals:compare}
\begin{enumerate}
  \item Primero compare las partes enteras: $3.2 > 2.9$.
  \item Mismas partes enteras: compara las décimas, luego las centésimas:
        $2.37 < 2.5$ porque $3$ décimas $< 5$ décimas.
  \item Escribir ambos con dos decimales ayuda: $2.5 = 2.50$, y
        $37 < 50$ centésimas.
\end{enumerate}
La trampa: $2.37$ tiene más dígitos que $2.5$, pero es
\emph{menor} --- ¡más dígitos no significa más grande!
\end{method}

\begin{example}\label{ex:g4:decimals:compare}
Orden $0.4$, $0.35$ y $0.09$: con centésimas, $0.40$, $0.35$,
$0.09$; entonces
\[
0.09 < 0.35 < 0.4 .
\]
\end{example}

\section{Ejercicios}

\begin{exercise}[$\star $]\label{exo:g4:decimals:1}
En cien \omterm{def:g3:shapes:shapes}{cuadrados}, ¿cuánto se sombrea si se sombrea? $5$ columnas;
$2$ columnas y $3$ células; $45$ células? Responde con fracciones.
\end{exercise}

\begin{exercise}[$\star $]\label{exo:g4:decimals:2}
escribe con un \omterm{def:g4:decimals:point}{coma decimal}:
$\frac{9}{10}$; \quad $\frac{27}{100}$; \quad $\frac{3}{100}$; \quad
$5 + \frac{6}{10}$; \quad $\frac{60}{100}$.
\end{exercise}

\begin{exercise}[$\star $]\label{exo:g4:decimals:3}
Escribe como un \omterm{def:g4:fractions:def}{fracción} (décimas o centésimas):
$0.3$; \quad $0.51$; \quad $0.07$; \quad $1.9$.
\end{exercise}

\begin{exercise}[$\star $]\label{exo:g4:decimals:4}
En $6.58$: ¿Qué significa el dígito? $5$ ¿contar? el dígito $8$? descomponer
el número como en \cref{def:g4:decimals:point}.
\end{exercise}

\begin{exercise}[$\star $]\label{exo:g4:decimals:5}
Colocar en una recta numérica graduada en décimas de $4$ a $5$:
$4.2$; \ $4.75$ (¿entre qué dos marcas?); \ $4.9$.
\end{exercise}

\begin{exercise}[$\star $]\label{exo:g4:decimals:6}
Copiar y completar con $<$, $>$ o $=$:
\[
0.6 \;?\; 0.60, \qquad
2.8 \;?\; 2.75, \qquad
0.09 \;?\; 0.1, \qquad
5.3 \;?\; 5.13 .
\]
\end{exercise}

\begin{exercise}[$\star $]\label{exo:g4:decimals:7}
Ordene de menor a mayor: $1.05$; \ $1.5$; \ $0.95$; \ $1.45$.
\end{exercise}

\begin{exercise}[$\star $]\label{exo:g4:decimals:8}
Escribir en escritura decimal: tres euros con cinco céntimos; doce euros
y cincuenta centavos; noventa y nueve centavos. Luego ordene los tres precios.
\end{exercise}

\begin{exercise}[$\star $]\label{exo:g4:decimals:9}
¿Qué números decimales con un dígito decimal se encuentran estrictamente entre
$3$ y $4$? Enumérelos todos. ¿Cuantos hay?
\end{exercise}

\begin{exercise}[$\star\star $]\label{exo:g4:decimals:10}
Tom dice: ``$0.25$ es más grande que $0.5$ porque $25$ es más grande que
$5$''. Muéstrale su error con la centena cuadrada (cuántas celdas
¿Cada tono?).
\end{exercise}

\begin{exercise}[$\star\star $]\label{exo:g4:decimals:11}
Se miden los tres saltos de un corredor $3.6$ metro, $3.58$ m y $3.65$ metro.
Ordena los saltos de más corto a más largo. ¿Cuáles dos se diferencian por
exactamente $\frac{5}{100}$ de un metro?
\end{exercise}
